\chapter{Contributing}
\label{chap:contributing}

\section{Scope of Contributions}

\actsim is intended to support open actuarial research and transparent
simulation workflows. Contributions are welcome and may include:

\begin{itemize}[leftmargin=*]
    \item Bug fixes.
    \item Support for additional probability distributions.
    \item New goodness-of-fit metrics.
    \item Additional validation tests.
    \item Documentation improvements (including this manual).
    \item Worked examples and case studies.
    \item Integration examples with other actuarial packages.
    \item Performance improvements to the simulation engine.
\end{itemize}

The package is distributed under the Apache License 2.0, and all
contributions are accepted under the same terms.

\section{Setting Up a Development Environment}

The recommended workflow uses an isolated virtual environment:

\begin{lstlisting}[style=bashstyle]
git clone https://github.com/jzhng105/actsim.git
cd actsim

python -m venv venv
source venv/bin/activate          # macOS / Linux
venv\Scripts\activate             # Windows

pip install -e .[dev]
\end{lstlisting}

Editable installation (\texttt{pip install -e}) ensures that any local
changes to the source files are picked up immediately, without
reinstallation.

\section{Running the Test Suite}

Tests are run using \texttt{pytest}:

\begin{lstlisting}[style=bashstyle]
pytest
\end{lstlisting}

When adding new functionality, please add corresponding tests under the
\texttt{tests/} directory. Tests should be deterministic and run
quickly.

\section{Code Style}

The code base follows standard Python conventions:

\begin{itemize}[leftmargin=*]
    \item PEP 8 formatting.
    \item Type hints on public method signatures where reasonable.
    \item Concise, descriptive names.
    \item Public methods documented with docstrings.
\end{itemize}

The repository's \texttt{black} configuration is the canonical source
for formatting choices. Running \texttt{black .} before committing is
recommended.

\section{Adding a New Distribution}

To add support for a new distribution:

\begin{enumerate}[leftmargin=*]
    \item Add the distribution to \texttt{actstats} (or confirm that an
          implementation already exists).
    \item Register the distribution in the
          \texttt{available\_distributions} dictionary in both
          \texttt{actfitter.py} and \texttt{actsimulator.py}.
    \item Add the distribution to the \texttt{config.yaml} file under
          the appropriate category (severity or frequency).
    \item Add a unit test that fits the new distribution to a known
          sample and verifies that the recovered parameters are close
          to the truth.
\end{enumerate}

\section{Pull Request Workflow}

\begin{enumerate}[leftmargin=*]
    \item Fork the repository on GitHub.
    \item Create a feature branch:
          \texttt{git checkout -b feature/short-description}.
    \item Make your changes and add tests.
    \item Run the test suite locally and ensure that it passes.
    \item Commit your changes with a descriptive message.
    \item Push the branch to your fork and open a pull request against
          the upstream repository's \texttt{main} branch.
\end{enumerate}

The maintainer will review the pull request and request changes if
necessary. Small, focused pull requests are easier to review than large
omnibus changes.

\section{Reporting Issues}

Bug reports and feature requests should be filed via the GitHub issue
tracker. Useful bug reports include:

\begin{itemize}[leftmargin=*]
    \item A short description of the unexpected behaviour.
    \item The version of \actsim, Python, and key dependencies.
    \item A minimal reproducible example.
    \item The actual and expected output.
\end{itemize}

\section{Discussion and Roadmap}

Open-ended discussion (design questions, methodology, roadmap) should
take place on the GitHub Discussions tab rather than the issue tracker.
The maintainer welcomes input on roadmap priorities, particularly from
practising actuaries who can help align development with real-world
needs.
